\chapter{Infinite Horizon Economies}
\section{Bubbles}
\subsection{Why is there no bubbles when $T$ is finite?}
\paragraph{[Answer]}
Intuitively, when $T$ is finite, the last asset holder must equal the price to its fundamental value,
which eliminates the possiblity of bubbles.
Interating backward, we have no bubble for each period. 
Formally, we can express asset prices according to 
\[
p_t = f_t + \mathbb{E}[m_{t,T} d_{t,T}]
\]
when $t = T$, the second term equals to 0.
Hence, there is no bubble when $T$ is finite. 